\section*{Bloque C. Casos con implicancias tarifarias y regulatorias del sector saneamiento}

\noindent\textit{Preguntas 66 a 90. Aplican la teoría del productor y de los costos a la interpretación de estructuras tarifarias, economías de escala y alcance, eficiencia productiva, productividad, agua no facturada y dilemas regulatorios de las empresas prestadoras de servicios de saneamiento (EPS).}

\q{Una EPS tiene una función de costo total $C(Q)=1000+2Q+0.1Q^2$, donde $Q$ se mide en miles de m$^3$. Si $Q=50$, el costo medio es:}
\begin{opts}
    \item 22
    \item 27
    \item 12
    \item 30
\end{opts}

\q{Con la misma función $C(Q)=1000+2Q+0.1Q^2$, el costo marginal en $Q=50$ es:}
\begin{opts}
    \item 2
    \item 7
    \item 22
    \item 12
\end{opts}

\q{Con $C(Q)=1000+2Q+0.1Q^2$, el nivel de producción que minimiza el costo medio es:}
\begin{opts}
    \item 50
    \item 70
    \item 100
    \item 150
\end{opts}

\q{Si el regulador fija una tarifa igual al costo marginal y la EPS produce $Q=50$, la pérdida que debe cubrirse con un subsidio o cargo fijo es:}
\begin{opts}
    \item 350
    \item 600
    \item 750
    \item 1350
\end{opts}

\q{Una empresa de saneamiento es un monopolio natural cuando:}
\begin{opts}
    \item el costo medio decrece y la función de costos es subaditiva.
    \item el costo medio es creciente en todo nivel de producción.
    \item el costo marginal es constante e igual al costo medio total.
    \item existen severas deseconomías de alcance entre los servicios.
\end{opts}

\q{Existen economías de alcance en una EPS que presta agua potable y alcantarillado cuando:}
\begin{opts}
    \item el costo medio del agua cae al ampliar solo su propio volumen.
    \item el costo marginal permanece constante en todo rango operativo.
    \item prestar ambos servicios juntos cuesta menos que por separado.
    \item los costos fijos de infraestructura resultan iguales a cero.
\end{opts}

\q{La regulación por precio tope RPI-X incentiva a la empresa a reducir costos porque:}
\begin{opts}
    \item retiene las ganancias de eficiencia dentro del período regulado.
    \item el regulador reembolsa automáticamente cualquier costo incurrido.
    \item puede sobrecapitalizar para ampliar la base tarifaria regulada.
    \item la tarifa anual se iguala siempre al costo medio observado.
\end{opts}

\q{El efecto Averch-Johnson, asociado a la regulación por tasa de retorno, tiende a generar:}
\begin{opts}
    \item sobrecapitalización para ampliar la base de capital regulada.
    \item subcapitalización extrema al no remunerar el capital invertido.
    \item competencia perfecta inmediata en el sector de saneamiento.
    \item eliminación total de los costos hundidos de infraestructura.
\end{opts}

\q{La regulación por costo de servicio (cost-plus) presenta el riesgo de:}
\begin{opts}
    \item incentivos sumamente fuertes para minimizar costos internos.
    \item tarifas fijas totalmente independientes del costo observado.
    \item competencia directa por comparación entre distintas firmas.
    \item incentivos débiles a la eficiencia y sesgo a inflar costos.
\end{opts}

\q{La competencia por comparación o yardstick competition consiste en:}
\begin{opts}
    \item fijar la tarifa uniforme exactamente igual al costo marginal.
    \item comparar el desempeño con empresas pares para fijar metas.
    \item regular exclusivamente la tasa de retorno sobre activos.
    \item permitir a la EPS cobrar libremente el precio de monopolio.
\end{opts}

\q{Si el precio de la energía eléctrica, insumo intensivo en el bombeo de agua, aumenta de forma permanente, en el largo plazo la EPS tenderá a:}
\begin{opts}
    \item sustituir energía por capital mediante equipos eficientes.
    \item sustituir capital físico por mayor uso intensivo de energía.
    \item sustituir mano de obra calificada por energía eléctrica.
    \item mantener rígida la misma proporción histórica de insumos.
\end{opts}

\q{Una EPS produce 1,000,000 de m$^3$ al año con 50 trabajadores. Su productividad media del trabajo es:}
\begin{opts}
    \item 10,000 m$^3$ por trabajador.
    \item 50,000 m$^3$ por trabajador.
    \item 5,000 m$^3$ por trabajador.
    \item 20,000 m$^3$ por trabajador.
\end{opts}

\q{Si el producto de una EPS crece 6\% y sus insumos crecen 4\%, la productividad total de los factores (residual de Solow) crece aproximadamente:}
\begin{opts}
    \item 10\%.
    \item 6\%.
    \item 2\%.
    \item 4\%.
\end{opts}

\q{Si el agua no facturada (ANF) es 25\% del agua producida y el costo de producción es S/ 1.50 por m$^3$, el costo efectivo por m$^3$ facturado es:}
\begin{opts}
    \item S/ 1.50.
    \item S/ 1.875.
    \item S/ 2.25.
    \item S/ 2.00.
\end{opts}

\q{Si la ANF se reduce de 25\% a 15\%, manteniendo el costo de producción en S/ 1.50 por m$^3$, el costo efectivo por m$^3$ facturado se reduce en aproximadamente:}
\begin{opts}
    \item S/ 0.12.
    \item S/ 0.24.
    \item S/ 0.35.
    \item S/ 0.50.
\end{opts}

\q{Un programa de reducción de fugas cuesta S/ 80,000 al año y permite ahorrar 50,000 m$^3$ de agua cuyo costo marginal de producción es S/ 2 por m$^3$. El beneficio neto anual del programa es:}
\begin{opts}
    \item S/ 0.
    \item S/ 10,000.
    \item S/ 20,000.
    \item S/ 30,000.
\end{opts}

\q{En una tarifa por bloques horarios o de punta, la lógica de eficiencia indica que:}
\begin{opts}
    \item se cobre una tarifa idéntica en todas las franjas horarias.
    \item se cobre más en hora punta por el mayor costo marginal de capacidad.
    \item se cobre menos en hora punta para estimular el consumo nocturno.
    \item se fije tarifa cero durante los períodos de máxima congestión.
\end{opts}

\q{Cuando una inversión en redes es hundida y específica, el regulador puede enfrentar un problema de:}
\begin{opts}
    \item compromiso creíble: riesgo de no reconocimiento tras invertir.
    \item competencia perfecta entre múltiples redes superpuestas.
    \item ausencia completa de costos fijos en la etapa de distribución.
    \item inexistencia de economías de escala en tuberías principales.
\end{opts}

\q{Si la base de activos regulados es S/ 1,000,000 y la depreciación lineal es 5\% anual, la depreciación anual es:}
\begin{opts}
    \item S/ 50,000.
    \item S/ 100,000.
    \item S/ 5,000.
    \item S/ 20,000.
\end{opts}

\q{Un costo común de S/ 100,000 se reparte entre agua y alcantarillado según sus costos separados: 60\% agua y 40\% alcantarillado. La parte asignada al agua es:}
\begin{opts}
    \item S/ 40,000.
    \item S/ 50,000.
    \item S/ 60,000.
    \item S/ 100,000.
\end{opts}

\q{Si el costo fijo anual de una EPS es S/ 1,000,000 y tiene 10,000 usuarios, el cargo fijo puro por usuario para cubrirlo es:}
\begin{opts}
    \item S/ 10.
    \item S/ 50.
    \item S/ 1,000.
    \item S/ 100.
\end{opts}

\q{Si el costo fijo que debe recuperarse es S/ 500,000 y el volumen facturado es 100,000 m$^3$, el cargo volumétrico uniforme que lo cubre es:}
\begin{opts}
    \item S/ 0.50 por m$^3$.
    \item S/ 1.00 por m$^3$.
    \item S/ 2.00 por m$^3$.
    \item S/ 5.00 por m$^3$.
\end{opts}

\q{Si dos EPS pequeñas se fusionan y existen economías de escala en el rango relevante, cabe esperar que la fusión:}
\begin{opts}
    \item baje el costo medio si las economías superan la coordinación.
    \item suba inevitablemente el costo medio desde la primera etapa.
    \item carezca de todo impacto sobre los costos medios unitarios.
    \item solo funcione cuando existan economías de alcance y no escala.
\end{opts}

\q{Si el regulador exige mejorar la continuidad y la presión del servicio, la teoría de costos sugiere que:}
\begin{opts}
    \item mayor calidad eleva costos y exige reconocimiento tarifario.
    \item la mejora técnica es gratuita y no altera los costos reales.
    \item reducir calidad siempre baja costos sin dañar el bienestar.
    \item la calidad operativa es ajena a la inversión y mantenimiento.
\end{opts}

\q{Si el costo marginal de conservar agua mediante reducción de fugas es S/ 1.50 por m$^3$ y el costo marginal de ampliar una nueva fuente es S/ 2.50 por m$^3$, la opción económicamente preferible es:}
\begin{opts}
    \item ampliar de inmediato la nueva fuente.
    \item ambas alternativas tienen costo igual.
    \item mantener las fugas sin intervención.
    \item priorizar la reducción de pérdidas.
\end{opts}
